% Created 2018-02-01 Thu 15:00
% Intended LaTeX compiler: pdflatex
\documentclass[11pt]{article}
\usepackage[utf8]{inputenc}
\usepackage[T1]{fontenc}
\usepackage{graphicx}
\usepackage{grffile}
\usepackage{longtable}
\usepackage{wrapfig}
\usepackage{rotating}
\usepackage[normalem]{ulem}
\usepackage{amsmath}
\usepackage{textcomp}
\usepackage{amssymb}
\usepackage{capt-of}
\usepackage{hyperref}
\author{Monib Ahmed}
\date{\today}
\title{EXAMPLE PYTHON CODE}
\hypersetup{
 pdfauthor={Monib Ahmed},
 pdftitle={EXAMPLE PYTHON CODE},
 pdfkeywords={},
 pdfsubject={},
 pdfcreator={Emacs 25.3.1 (Org mode 9.1.5)}, 
 pdflang={English}}
\begin{document}

\maketitle
\tableofcontents


\section{Hello World Example}
\label{sec:orgcd5ddb5}
This is an hello world example to include some images.

\begin{verbatim}
import matplotlib.pyplot as plt
import numpy as np

# Data for plotting
t = np.arange(0.0, 2.0, 0.01)
s = 1 + np.sin(2 * np.pi * t)

# Note that using plt.subplots below is equivalent to using
# fig = plt.figure and then ax = fig.add_subplot(111)
fig, ax = plt.subplots()
ax.plot(t, s)

ax.set(xlabel='time (s)', ylabel='voltage (mV)',
       title='About as simple as it gets, folks')
ax.grid()

fig.savefig("test.png")
plt.show()
\end{verbatim}
\end{document}
